Since the COVID-19 pandemic, vaccination rates have fallen across a wide category of diseases. 
The share of kindergartners not vaccinated against measles, mumps, and urbella increased from 4 
percent in 2020 to 6 percent in 2025; polio and hepatitis-B vaccination rates fell similarly
\citep{cdc_schoolvaxview_data}.
Adult vaccination rates have fallen as well, with flu vaccination rates falling each year 
since 2020 \citep{cdc_flu_vaccinations_weekly}.
These declining vaccination rates are concerning given the health benefits of vaccination 
and widespread scientific evidence supporting the safety and effectiveness of vaccination.

Many interventions have been tested to address declining vaccination rates, but with limited 
success, especially for vaccine hesitant people. 
For example, nudges and reminders can raise flu vaccination rates by 1-3 percentage points 
\citep{milkman2022680,patel2023randomized}, but reminders had mixed effects among the hesitant 
\citep{chang2023reminders,chang2026impact} Incentives raised vaccination rates COVID booster 
rates by A-B percentage points, among people who had already received a primary shot 
\citep{campos2024incentives}. Incentives in the US among vaccine hesitant Americans produced 
much smaller increases \citep{chang2023reminders, chang2026impact}.

The limited success of nudges, reminders, and incentives indicates that Increasing vaccination rates among vaccine hesitant people may require changing their beliefs about the benefits and costs of vaccination, 
although doing so in practice has proven challenging. Addressing beliefs via information provision has yielded mixed results: dispelling myths did not increase flu or pediatric vaccination rates 
\citep{nyhan2014effective, nyhan2015does}, 
while highlighting harms of pediatric illness raised support for pediatric vaccinations \citep{horne2015countering}.
Endorsements from trusted messengers may help \citep{larsen2023counter,alsan2024experimental}, albeit with modest effect sizes.  

In this paper, we test whether providing scientific evidence about flu vaccine side effects to vaccine hesitant people can shift beliefs about their own risks from the flu vaccine. 
Scientific information provides part of the basis for the widespread recommendation to evaluate---establishing the safety of vaccines---so it is natural to communicate. However it may not influence 
beliefs if vaccine hesitant participants do not trust the information source or do not believe it is relevant to them. 

We focus on the flu vaccine because flu remains an important seasonal disease, killing tens of thousands of Americans annually and infecting tens of millions. 
We focus on side effect beliefs because prior work has shown that flu vaccine hesitancy is strongly associated with pessimistic beliefs about sides \citep{chapman1999predictors,tsutsui2012economic}.
Quantitative belief elicitations show that vaccine hesitant people believe the flu vaccine side effect rate is 7-15 percentage points on average, 
much higher than clinical trial reports \citep{sacks2025preferences,fda_flulaval_package_insert,fda_fluzone_quadrivalent_package_insert}. 
Thus there is scope to move side effect beliefs with scientific information. Flu vaccine hesitancy is also associated with more negative views about 
flu vaccine effectiveness and side effect risk, but those views are nonetheless consistent with CDC estimates \citep{sacks2025preferences}.

We use an online experiment with 3,525 participants to test three forms of information provision. We ask three research questions. 
First, can scientific information provision to vaccine hesitant people 
change their beliefs about side effect risks? Second, motivated by the possibility that vaccine hesitancy is associated with mistrust of the pharmaceutical industry. 

We use an online experiment with 3,525 vaccine hesitant participants to estimate the effects of flu vaccine side effect information provision. 
In our control condition, we inform participants that vaccine manufacturers use placebo-controlled clinical trials to establish the effectiveness and safety of vaccines, 
and we tell them that the severe adverse event rate in the placebo arm of a flu vaccine trial was 3 percent.  In our three treatment conditions, 
we provide information about the low adverse event rate in the vaccine arms. One condition describes results from the industry-sponsored trial. 
A second condition reports low adverse event rate in an academic study, and a third emphasizes the representativeness of participants in the academic study to current participants. 
We describe this as scientific information provision because we describe the setup and results of research studies. 
We measure beliefs prior to the intervention and, in a modified form, after the intervention. We measure beliefs both about the trial adverse event rate and about personal side effect risk, our primary outcome. 
These belief measures are correlated with proxies for vaccine hesitancy (such as trust in doctors or willingness to follow government recommendations), as well as with prior experiences with the flu and COVID vaccines. 

We have four main results. First, scientific information provision can reduce beliefs about side effect risk among the vaccine hesitant. 
In clinical trials, the severe adverse event rate in the vaccine and placebo arms are within about 1 percentage points of each other.
In our control group, vaccine hesitant people believe this difference is 15 percentage points on average. In our industry information arm, this belief was 3 percentage points lower.  
Our second result is that funding source—academic or industry—does not affect belief updating. 
Third, our arm highlighting representativeness did not increase perceived relevance of information. 
Fourth, scientific information provision can increase intentions to vaccinate, which grew by 3 percentage points (45\% of control mean) in the academic arm. 
This increased intention did not translate into a statistically significant increase in vaccinations (measured two weeks later in a follow-up survey), 
although we lacked the power to detect plausible increases in vaccinations. This lack of power for vaccine take-up is an important limitation of our study, 
as is our reliance on subjective measures. 

Several results confirm that our information led to belief updating, and rule out other explanations such as experimenter demand or expressive answering. 
First, we find larger effects for participants with more negative prior beliefs, as well as for participants without prior vaccine experience (whose beliefs should be easier to influence). 
Second, we incentivized participants to report beliefs about trial outcomes. These incentivized beliefs are correlated with reported personal risk. 
Third, our treatment effects on beliefs and intentions are moderated by trial perceptions of trust and relevance. 
For participants who found the provided information highly trustworthy and relevant, the effects were large. 
For participants who found it untrustworthy or irrelevant, the effects on beliefs were small and the effects on intentions were negative, 
a backlash effect similar to that documented by \citet{nyhan2014effective}. This pattern of responses and heterogeneous effects is consistent with belief reporting and 
updating but not experimenter demand or expressive answering.  

Our results point to a potentially promising class of interventions. 
Providing scientific information about side effects appears effective at moving beliefs and flu vaccine intentions among the vaccine hesitant. A
lthough vaccine hesitancy is characterized by mistrust of doctors and government recommendations, it is nonetheless possible to move some participants’ beliefs with scientific information. 
As the updating happened among people who trusted the trial information, our results suggest that side effect information provision would be especially effective if delivered via a trusted messenger, 
as prior work has found. Prior work such as \citet{alsan2024experimental,larsen2023counter} has shown the importance of the messenger; our contribution is to find a promising message. 



